% ============================================================
% frontmatter/abstract_cn.tex —— 中文摘要 + 关键词
% 撰写要求（两刊对照，投稿时按目标刊取严）：
%   - 《系统工程与电子技术》：150～250 字，关键词 4～8 个
%   - 《信息与控制》：≤ 500 字，关键词 3～6 个（建议用 EI 主题词）
%   - 取严写法：摘要 150～250 字，关键词 4～6 个
%   - 四要素：目的 — 方法 — 结果 — 结论；避免公式/图表/参考文献编号/非常用缩略语
%   - 关键词按“研究对象—方法—性质”顺序排列
% ============================================================

\vspace{1em}

\noindent{\heiti\zihao{5}摘要：}%
    面向任务持续到达、机巢分散部署的广域常态化无人机作业场景，作业过程中常伴随无人机故障、紧急任务等突发扰动。现有动态多无人机任务分配研究多采用事件驱动重分配范式，但普遍隐含“到达任务应尽量全部分配”的假设，较少在机队能力不足时将任务选择作为显式决策；起降结构多简化为单基地或固定返航，终点机巢鲜少进入联合优化。当吞吐受限、机巢分散与地形约束并存时，需要在滚动时域内同时决定执行哪些任务、由何机按何顺序访问以及降落于何巢。
    本文将该问题表述为滚动时域下的选择性多机巢多无人机任务分配：在硬航程与结构约束下优先保障可完成任务的覆盖，覆盖等价时再压缩含返航的航程；允许任务暂缓，并将迟到作为软约束计入代价而非淘汰条件。方法采用两阶段分解。第一阶段在能力上界约束下从动态任务池中确定本周期承诺集，以控制后续组合规模并抑制任务被无限期推迟；第二阶段仅在已选集上求解多机归属、访问顺序与终点机巢，通过搜索与约束修复生成可行航次，使覆盖优先与航程优化在可行域内分层落实。滚动过程吸收新任务到达、平台故障与机巢状态变化，并对受影响子集实施部分重分配，以保持分配方案与动态环境一致。

    （占位）在数字高程模型驱动的多机巢场景及多随机种子实验中，所提框架在可完成任务集上保持高完成率，相对简化距离模型与构造性策略在单位完成航程与时效指标上具有优势，并在故障与紧急任务注入下维持任务守恒与协议一致的重规划行为。结果表明，将任务选择、多机巢终点决策与地形感知代价纳入统一滚动分配框架，可在动态资源约束与突发扰动下实现覆盖优先、成本可控的多机协同调度。

\par\vspace{0.8em}

\noindent{\heiti\zihao{5}关键词：}%
\placeholder{关键词一；关键词二；关键词三；关键词四；关键词五（4～6 个为宜）}

\vspace{1.2em}
